
\documentclass{article}
\usepackage{colortbl}
\usepackage{makecell}
\usepackage{multirow}
\usepackage{supertabular}

\begin{document}

\newcounter{utterance}

\centering \large Interaction Transcript for game `matchit\_ascii', experiment `similar\_grid\_2', episode 0 with claude{-}opus{-}4{-}8{-}azure.
\vspace{24pt}

{ \footnotesize  \setcounter{utterance}{1}
\setlength{\tabcolsep}{0pt}
\begin{supertabular}{c@{$\;$}|p{.15\linewidth}@{}p{.15\linewidth}p{.15\linewidth}p{.15\linewidth}p{.15\linewidth}p{.15\linewidth}}
    \# & $\;$A & \multicolumn{4}{c}{Game Master} & $\;\:$B\\
    \hline

    \theutterance \stepcounter{utterance}  
    & & \multicolumn{4}{p{0.6\linewidth}}{
        \cellcolor[rgb]{0.9,0.9,0.9}{
            \makecell[{{p{\linewidth}}}]{
                \texttt{\tiny{[P1$\langle$GM]}}
                \texttt{U neemt deel aan een gezamenlijk raadspel.} \\
\texttt{Het doel is om erachter te komen of dit raster en een andere afbeelding die alleen ik kan zien, hetzelfde zijn.} \\
\texttt{Beschrijf eerst uw raster. Daarna geef ik mijn beschrijving en kunnen we elkaar vragen stellen over de rasters om erachter te komen of ze hetzelfde zijn. Begin nu uw korte rasterbeschrijving met "BESCHRIJVING:" gevolgd door de beschrijving. Voeg niets anders toe. Dit is uw raster:} \\
\texttt{▢ X X X X} \\
\texttt{▢ X X X X} \\
\texttt{▢ X X X X} \\
\texttt{▢ X X X X} \\
\texttt{▢ ▢ ▢ ▢ ▢} \\
            }
        }
    }
    & \\ \\

    \theutterance \stepcounter{utterance}  
    & \multicolumn{4}{p{0.6\linewidth}}{
        \cellcolor[rgb]{0.9,0.9,0.9}{
            \makecell[{{p{\linewidth}}}]{
                \texttt{\tiny{[P1$\rangle$GM]}}
                \texttt{BESCHRIJVING: Een raster van 5 rijen en 5 kolommen. De eerste kolom bestaat volledig uit lege vakjes (▢). De onderste rij bestaat volledig uit lege vakjes (▢). Alle overige vakjes (de rechterbovenhoek, een blok van 4 rijen bij 4 kolommen) zijn gevuld met X.} \\
            }
        }
    }
    & & \\ \\

    \theutterance \stepcounter{utterance}  
    & & & \multicolumn{2}{p{0.3\linewidth}}{
        \cellcolor[rgb]{0.9,0.9,0.9}{
            \makecell[{{p{\linewidth}}}]{
                \texttt{\tiny{[GM$|$GM]}}
                \texttt{continue} \\
            }
        }
    }
    & & \\ \\

    \theutterance \stepcounter{utterance}  
    & & \multicolumn{4}{p{0.6\linewidth}}{
        \cellcolor[rgb]{0.9,0.9,0.9}{
            \makecell[{{p{\linewidth}}}]{
                \texttt{\tiny{[GM$\rangle$P2]}}
                \texttt{U neemt deel aan een gezamenlijk raadspel.} \\
\texttt{Het doel is om erachter te komen of dit raster en een andere afbeelding die alleen ik kan zien, hetzelfde zijn.} \\
\texttt{Beschrijf eerst uw raster. Daarna geef ik mijn beschrijving en kunnen we elkaar vragen stellen over de rasters om erachter te komen of ze hetzelfde zijn. Begin nu uw korte rasterbeschrijving met "BESCHRIJVING:" gevolgd door de beschrijving. Voeg niets anders toe. Dit is uw raster:} \\
\texttt{X X X X X} \\
\texttt{▢ X X X ▢} \\
\texttt{▢ X X X X} \\
\texttt{▢ X X X X} \\
\texttt{▢ ▢ ▢ ▢ ▢} \\
            }
        }
    }
    & \\ \\

    \theutterance \stepcounter{utterance}  
    & & & \multicolumn{4}{p{0.6\linewidth}}{
        \cellcolor[rgb]{0.9,0.9,0.9}{
            \makecell[{{p{\linewidth}}}]{
                \texttt{\tiny{[GM$\langle$P2]}}
                \texttt{BESCHRIJVING: Een 5x5 raster. Rij 1: alle vijf gevuld. Rij 2: leeg, gevuld, gevuld, gevuld, leeg. Rij 3: leeg, dan vier gevuld. Rij 4: leeg, dan vier gevuld. Rij 5: alle vijf leeg.} \\
            }
        }
    }
     \\ \\

    \theutterance \stepcounter{utterance}  
    & & & \multicolumn{2}{p{0.3\linewidth}}{
        \cellcolor[rgb]{0.9,0.9,0.9}{
            \makecell[{{p{\linewidth}}}]{
                \texttt{\tiny{[GM$|$GM]}}
                \texttt{continue} \\
            }
        }
    }
    & & \\ \\

    \theutterance \stepcounter{utterance}  
    & & \multicolumn{4}{p{0.6\linewidth}}{
        \cellcolor[rgb]{0.9,0.9,0.9}{
            \makecell[{{p{\linewidth}}}]{
                \texttt{\tiny{[GM$\rangle$P2]}}
                \texttt{Dit is mijn BESCHRIJVINGBESCHRIJVING: Een raster van 5 rijen en 5 kolommen. De eerste kolom bestaat volledig uit lege vakjes (▢). De onderste rij bestaat volledig uit lege vakjes (▢). Alle overige vakjes (de rechterbovenhoek, een blok van 4 rijen bij 4 kolommen) zijn gevuld met X.} \\
\texttt{Stel nu een vraag om nieuwe aspecten van mijn raster te ontdekken die anders kunnen zijn dan jouw raster. Begin met "VRAAG:" en voeg niets anders toe.} \\
            }
        }
    }
    & \\ \\

    \theutterance \stepcounter{utterance}  
    & & & \multicolumn{4}{p{0.6\linewidth}}{
        \cellcolor[rgb]{0.9,0.9,0.9}{
            \makecell[{{p{\linewidth}}}]{
                \texttt{\tiny{[GM$\langle$P2]}}
                \texttt{VRAAG: In jouw raster, is het eerste vakje van de bovenste rij (rij 1, kolom 1) gevuld met X of leeg?} \\
            }
        }
    }
     \\ \\

    \theutterance \stepcounter{utterance}  
    & & & \multicolumn{2}{p{0.3\linewidth}}{
        \cellcolor[rgb]{0.9,0.9,0.9}{
            \makecell[{{p{\linewidth}}}]{
                \texttt{\tiny{[GM$|$GM]}}
                \texttt{continue} \\
            }
        }
    }
    & & \\ \\

    \theutterance \stepcounter{utterance}  
    & & \multicolumn{4}{p{0.6\linewidth}}{
        \cellcolor[rgb]{0.9,0.9,0.9}{
            \makecell[{{p{\linewidth}}}]{
                \texttt{\tiny{[P1$\langle$GM]}}
                \texttt{Dit is mijn BESCHRIJVINGBESCHRIJVING: Een 5x5 raster. Rij 1: alle vijf gevuld. Rij 2: leeg, gevuld, gevuld, gevuld, leeg. Rij 3: leeg, dan vier gevuld. Rij 4: leeg, dan vier gevuld. Rij 5: alle vijf leeg.} \\
\texttt{VRAAG: In jouw raster, is het eerste vakje van de bovenste rij (rij 1, kolom 1) gevuld met X of leeg?Begin je antwoord met "ANTWOORD:" en voeg niets anders toe.} \\
            }
        }
    }
    & \\ \\

    \theutterance \stepcounter{utterance}  
    & \multicolumn{4}{p{0.6\linewidth}}{
        \cellcolor[rgb]{0.9,0.9,0.9}{
            \makecell[{{p{\linewidth}}}]{
                \texttt{\tiny{[P1$\rangle$GM]}}
                \texttt{ANTWOORD: In mijn raster is het eerste vakje van de bovenste rij (rij 1, kolom 1) leeg (▢).} \\
            }
        }
    }
    & & \\ \\

    \theutterance \stepcounter{utterance}  
    & & & \multicolumn{2}{p{0.3\linewidth}}{
        \cellcolor[rgb]{0.9,0.9,0.9}{
            \makecell[{{p{\linewidth}}}]{
                \texttt{\tiny{[GM$|$GM]}}
                \texttt{continue} \\
            }
        }
    }
    & & \\ \\

    \theutterance \stepcounter{utterance}  
    & & \multicolumn{4}{p{0.6\linewidth}}{
        \cellcolor[rgb]{0.9,0.9,0.9}{
            \makecell[{{p{\linewidth}}}]{
                \texttt{\tiny{[P1$\langle$GM]}}
                \texttt{Stel nu een vraag om nieuwe aspecten van mijn raster te ontdekken die anders kunnen zijn dan jouw raster. Begin met "VRAAG:" en voeg niets anders toe.} \\
            }
        }
    }
    & \\ \\

    \theutterance \stepcounter{utterance}  
    & \multicolumn{4}{p{0.6\linewidth}}{
        \cellcolor[rgb]{0.9,0.9,0.9}{
            \makecell[{{p{\linewidth}}}]{
                \texttt{\tiny{[P1$\rangle$GM]}}
                \texttt{VRAAG: In jouw raster, is het laatste vakje van rij 2 (rij 2, kolom 5) gevuld of leeg?} \\
            }
        }
    }
    & & \\ \\

    \theutterance \stepcounter{utterance}  
    & & & \multicolumn{2}{p{0.3\linewidth}}{
        \cellcolor[rgb]{0.9,0.9,0.9}{
            \makecell[{{p{\linewidth}}}]{
                \texttt{\tiny{[GM$|$GM]}}
                \texttt{continue} \\
            }
        }
    }
    & & \\ \\

    \theutterance \stepcounter{utterance}  
    & & \multicolumn{4}{p{0.6\linewidth}}{
        \cellcolor[rgb]{0.9,0.9,0.9}{
            \makecell[{{p{\linewidth}}}]{
                \texttt{\tiny{[GM$\rangle$P2]}}
                \texttt{ANTWOORD: In mijn raster is het eerste vakje van de bovenste rij (rij 1, kolom 1) leeg (▢).} \\
\texttt{VRAAG: In jouw raster, is het laatste vakje van rij 2 (rij 2, kolom 5) gevuld of leeg?Begin je antwoord met "ANTWOORD:" en voeg niets anders toe.} \\
            }
        }
    }
    & \\ \\

    \theutterance \stepcounter{utterance}  
    & & & \multicolumn{4}{p{0.6\linewidth}}{
        \cellcolor[rgb]{0.9,0.9,0.9}{
            \makecell[{{p{\linewidth}}}]{
                \texttt{\tiny{[GM$\langle$P2]}}
                \texttt{ANTWOORD: In mijn raster is het laatste vakje van rij 2 (rij 2, kolom 5) leeg (▢). Bovendien is bij mij rij 1, kolom 1 gevuld (X), terwijl die bij jou leeg is. De rasters zijn dus verschillend.} \\
            }
        }
    }
     \\ \\

    \theutterance \stepcounter{utterance}  
    & & & \multicolumn{2}{p{0.3\linewidth}}{
        \cellcolor[rgb]{0.9,0.9,0.9}{
            \makecell[{{p{\linewidth}}}]{
                \texttt{\tiny{[GM$|$GM]}}
                \texttt{continue} \\
            }
        }
    }
    & & \\ \\

    \theutterance \stepcounter{utterance}  
    & & \multicolumn{4}{p{0.6\linewidth}}{
        \cellcolor[rgb]{0.9,0.9,0.9}{
            \makecell[{{p{\linewidth}}}]{
                \texttt{\tiny{[GM$\rangle$P2]}}
                \texttt{Stel nu een vraag om nieuwe aspecten van mijn raster te ontdekken die anders kunnen zijn dan jouw raster. Begin met "VRAAG:" en voeg niets anders toe.} \\
            }
        }
    }
    & \\ \\

    \theutterance \stepcounter{utterance}  
    & & & \multicolumn{4}{p{0.6\linewidth}}{
        \cellcolor[rgb]{0.9,0.9,0.9}{
            \makecell[{{p{\linewidth}}}]{
                \texttt{\tiny{[GM$\langle$P2]}}
                \texttt{VRAAG: In jouw raster, is het laatste vakje van rij 3 (rij 3, kolom 5) gevuld of leeg?} \\
            }
        }
    }
     \\ \\

    \theutterance \stepcounter{utterance}  
    & & & \multicolumn{2}{p{0.3\linewidth}}{
        \cellcolor[rgb]{0.9,0.9,0.9}{
            \makecell[{{p{\linewidth}}}]{
                \texttt{\tiny{[GM$|$GM]}}
                \texttt{continue} \\
            }
        }
    }
    & & \\ \\

    \theutterance \stepcounter{utterance}  
    & & \multicolumn{4}{p{0.6\linewidth}}{
        \cellcolor[rgb]{0.9,0.9,0.9}{
            \makecell[{{p{\linewidth}}}]{
                \texttt{\tiny{[P1$\langle$GM]}}
                \texttt{ANTWOORD: In mijn raster is het laatste vakje van rij 2 (rij 2, kolom 5) leeg (▢). Bovendien is bij mij rij 1, kolom 1 gevuld (X), terwijl die bij jou leeg is. De rasters zijn dus verschillend.} \\
\texttt{VRAAG: In jouw raster, is het laatste vakje van rij 3 (rij 3, kolom 5) gevuld of leeg?Begin je antwoord met "ANTWOORD:" en voeg niets anders toe.} \\
            }
        }
    }
    & \\ \\

    \theutterance \stepcounter{utterance}  
    & \multicolumn{4}{p{0.6\linewidth}}{
        \cellcolor[rgb]{0.9,0.9,0.9}{
            \makecell[{{p{\linewidth}}}]{
                \texttt{\tiny{[P1$\rangle$GM]}}
                \texttt{ANTWOORD: In mijn raster is het laatste vakje van rij 3 (rij 3, kolom 5) gevuld met X. De rasters zijn inderdaad verschillend, want in mijn raster is rij 1, kolom 1 leeg terwijl die bij jou gevuld is, en in jouw raster is rij 2, kolom 5 leeg terwijl die bij mij gevuld is.} \\
            }
        }
    }
    & & \\ \\

    \theutterance \stepcounter{utterance}  
    & & & \multicolumn{2}{p{0.3\linewidth}}{
        \cellcolor[rgb]{0.9,0.9,0.9}{
            \makecell[{{p{\linewidth}}}]{
                \texttt{\tiny{[GM$|$GM]}}
                \texttt{continue} \\
            }
        }
    }
    & & \\ \\

    \theutterance \stepcounter{utterance}  
    & & \multicolumn{4}{p{0.6\linewidth}}{
        \cellcolor[rgb]{0.9,0.9,0.9}{
            \makecell[{{p{\linewidth}}}]{
                \texttt{\tiny{[P1$\langle$GM]}}
                \texttt{Stel nu een vraag om nieuwe aspecten van mijn raster te ontdekken die anders kunnen zijn dan jouw raster. Begin met "VRAAG:" en voeg niets anders toe.} \\
            }
        }
    }
    & \\ \\

    \theutterance \stepcounter{utterance}  
    & \multicolumn{4}{p{0.6\linewidth}}{
        \cellcolor[rgb]{0.9,0.9,0.9}{
            \makecell[{{p{\linewidth}}}]{
                \texttt{\tiny{[P1$\rangle$GM]}}
                \texttt{VRAAG: In jouw raster, is het eerste vakje van rij 5 (rij 5, kolom 1) gevuld of leeg?} \\
            }
        }
    }
    & & \\ \\

    \theutterance \stepcounter{utterance}  
    & & & \multicolumn{2}{p{0.3\linewidth}}{
        \cellcolor[rgb]{0.9,0.9,0.9}{
            \makecell[{{p{\linewidth}}}]{
                \texttt{\tiny{[GM$|$GM]}}
                \texttt{continue} \\
            }
        }
    }
    & & \\ \\

    \theutterance \stepcounter{utterance}  
    & & \multicolumn{4}{p{0.6\linewidth}}{
        \cellcolor[rgb]{0.9,0.9,0.9}{
            \makecell[{{p{\linewidth}}}]{
                \texttt{\tiny{[GM$\rangle$P2]}}
                \texttt{ANTWOORD: In mijn raster is het laatste vakje van rij 3 (rij 3, kolom 5) gevuld met X. De rasters zijn inderdaad verschillend, want in mijn raster is rij 1, kolom 1 leeg terwijl die bij jou gevuld is, en in jouw raster is rij 2, kolom 5 leeg terwijl die bij mij gevuld is.} \\
\texttt{VRAAG: In jouw raster, is het eerste vakje van rij 5 (rij 5, kolom 1) gevuld of leeg?Begin je antwoord met "ANTWOORD:" en voeg niets anders toe.} \\
            }
        }
    }
    & \\ \\

    \theutterance \stepcounter{utterance}  
    & & & \multicolumn{4}{p{0.6\linewidth}}{
        \cellcolor[rgb]{0.9,0.9,0.9}{
            \makecell[{{p{\linewidth}}}]{
                \texttt{\tiny{[GM$\langle$P2]}}
                \texttt{ANTWOORD: In mijn raster is het eerste vakje van rij 5 (rij 5, kolom 1) leeg (▢). Dat is hetzelfde als bij jou. Maar we hebben al vastgesteld dat de rasters verschillen: bij mij is rij 1, kolom 1 gevuld (bij jou leeg) en bij jou is rij 2, kolom 5 leeg (bij mij gevuld). De rasters zijn dus verschillend.} \\
            }
        }
    }
     \\ \\

    \theutterance \stepcounter{utterance}  
    & & & \multicolumn{2}{p{0.3\linewidth}}{
        \cellcolor[rgb]{0.9,0.9,0.9}{
            \makecell[{{p{\linewidth}}}]{
                \texttt{\tiny{[GM$|$GM]}}
                \texttt{continue} \\
            }
        }
    }
    & & \\ \\

    \theutterance \stepcounter{utterance}  
    & & \multicolumn{4}{p{0.6\linewidth}}{
        \cellcolor[rgb]{0.9,0.9,0.9}{
            \makecell[{{p{\linewidth}}}]{
                \texttt{\tiny{[GM$\rangle$P2]}}
                \texttt{Stel nu een vraag om nieuwe aspecten van mijn raster te ontdekken die anders kunnen zijn dan jouw raster. Begin met "VRAAG:" en voeg niets anders toe.} \\
            }
        }
    }
    & \\ \\

    \theutterance \stepcounter{utterance}  
    & & & \multicolumn{4}{p{0.6\linewidth}}{
        \cellcolor[rgb]{0.9,0.9,0.9}{
            \makecell[{{p{\linewidth}}}]{
                \texttt{\tiny{[GM$\langle$P2]}}
                \texttt{VRAAG: In jouw raster, is het laatste vakje van rij 4 (rij 4, kolom 5) gevuld of leeg?} \\
            }
        }
    }
     \\ \\

    \theutterance \stepcounter{utterance}  
    & & & \multicolumn{2}{p{0.3\linewidth}}{
        \cellcolor[rgb]{0.9,0.9,0.9}{
            \makecell[{{p{\linewidth}}}]{
                \texttt{\tiny{[GM$|$GM]}}
                \texttt{continue} \\
            }
        }
    }
    & & \\ \\

    \theutterance \stepcounter{utterance}  
    & & \multicolumn{4}{p{0.6\linewidth}}{
        \cellcolor[rgb]{0.9,0.9,0.9}{
            \makecell[{{p{\linewidth}}}]{
                \texttt{\tiny{[P1$\langle$GM]}}
                \texttt{ANTWOORD: In mijn raster is het eerste vakje van rij 5 (rij 5, kolom 1) leeg (▢). Dat is hetzelfde als bij jou. Maar we hebben al vastgesteld dat de rasters verschillen: bij mij is rij 1, kolom 1 gevuld (bij jou leeg) en bij jou is rij 2, kolom 5 leeg (bij mij gevuld). De rasters zijn dus verschillend.} \\
\texttt{VRAAG: In jouw raster, is het laatste vakje van rij 4 (rij 4, kolom 5) gevuld of leeg?Begin je antwoord met "ANTWOORD:" en voeg niets anders toe.} \\
            }
        }
    }
    & \\ \\

    \theutterance \stepcounter{utterance}  
    & \multicolumn{4}{p{0.6\linewidth}}{
        \cellcolor[rgb]{0.9,0.9,0.9}{
            \makecell[{{p{\linewidth}}}]{
                \texttt{\tiny{[P1$\rangle$GM]}}
                \texttt{ANTWOORD: In mijn raster is het laatste vakje van rij 4 (rij 4, kolom 5) gevuld met X. We hebben al vastgesteld dat de rasters verschillend zijn, want in mijn raster is rij 1, kolom 1 leeg terwijl die bij jou gevuld is, en in jouw raster is rij 2, kolom 5 leeg terwijl die bij mij gevuld is.} \\
            }
        }
    }
    & & \\ \\

    \theutterance \stepcounter{utterance}  
    & & & \multicolumn{2}{p{0.3\linewidth}}{
        \cellcolor[rgb]{0.9,0.9,0.9}{
            \makecell[{{p{\linewidth}}}]{
                \texttt{\tiny{[GM$|$GM]}}
                \texttt{continue} \\
            }
        }
    }
    & & \\ \\

    \theutterance \stepcounter{utterance}  
    & & \multicolumn{4}{p{0.6\linewidth}}{
        \cellcolor[rgb]{0.9,0.9,0.9}{
            \makecell[{{p{\linewidth}}}]{
                \texttt{\tiny{[P1$\langle$GM]}}
                \texttt{Stel nu een vraag om nieuwe aspecten van mijn raster te ontdekken die anders kunnen zijn dan jouw raster. Begin met "VRAAG:" en voeg niets anders toe.} \\
            }
        }
    }
    & \\ \\

    \theutterance \stepcounter{utterance}  
    & \multicolumn{4}{p{0.6\linewidth}}{
        \cellcolor[rgb]{0.9,0.9,0.9}{
            \makecell[{{p{\linewidth}}}]{
                \texttt{\tiny{[P1$\rangle$GM]}}
                \texttt{VRAAG: In jouw raster, is het eerste vakje van rij 3 (rij 3, kolom 1) gevuld of leeg?} \\
            }
        }
    }
    & & \\ \\

    \theutterance \stepcounter{utterance}  
    & & & \multicolumn{2}{p{0.3\linewidth}}{
        \cellcolor[rgb]{0.9,0.9,0.9}{
            \makecell[{{p{\linewidth}}}]{
                \texttt{\tiny{[GM$|$GM]}}
                \texttt{continue} \\
            }
        }
    }
    & & \\ \\

    \theutterance \stepcounter{utterance}  
    & & \multicolumn{4}{p{0.6\linewidth}}{
        \cellcolor[rgb]{0.9,0.9,0.9}{
            \makecell[{{p{\linewidth}}}]{
                \texttt{\tiny{[GM$\rangle$P2]}}
                \texttt{ANTWOORD: In mijn raster is het laatste vakje van rij 4 (rij 4, kolom 5) gevuld met X. We hebben al vastgesteld dat de rasters verschillend zijn, want in mijn raster is rij 1, kolom 1 leeg terwijl die bij jou gevuld is, en in jouw raster is rij 2, kolom 5 leeg terwijl die bij mij gevuld is.} \\
\texttt{VRAAG: In jouw raster, is het eerste vakje van rij 3 (rij 3, kolom 1) gevuld of leeg?Begin je antwoord met "ANTWOORD:" en voeg niets anders toe.} \\
            }
        }
    }
    & \\ \\

    \theutterance \stepcounter{utterance}  
    & & & \multicolumn{4}{p{0.6\linewidth}}{
        \cellcolor[rgb]{0.9,0.9,0.9}{
            \makecell[{{p{\linewidth}}}]{
                \texttt{\tiny{[GM$\langle$P2]}}
                \texttt{ANTWOORD: In mijn raster is het eerste vakje van rij 3 (rij 3, kolom 1) leeg (▢). Dat is hetzelfde als bij jou. Maar de rasters zijn nog steeds verschillend, zoals eerder vastgesteld: bij mij is rij 1, kolom 1 gevuld (bij jou leeg) en bij jou is rij 2, kolom 5 leeg (bij mij gevuld).} \\
            }
        }
    }
     \\ \\

    \theutterance \stepcounter{utterance}  
    & & & \multicolumn{2}{p{0.3\linewidth}}{
        \cellcolor[rgb]{0.9,0.9,0.9}{
            \makecell[{{p{\linewidth}}}]{
                \texttt{\tiny{[GM$|$GM]}}
                \texttt{continue} \\
            }
        }
    }
    & & \\ \\

    \theutterance \stepcounter{utterance}  
    & & \multicolumn{4}{p{0.6\linewidth}}{
        \cellcolor[rgb]{0.9,0.9,0.9}{
            \makecell[{{p{\linewidth}}}]{
                \texttt{\tiny{[GM$\rangle$P2]}}
                \texttt{Neem nu een beslissing. Wat denk je: zijn jouw raster en het andere beschreven raster hetzelfde? Schrijf "BESLISSING: hetzelfde raster" als je denkt dat het hetzelfde raster is of "BESLISSING: verschillende rasters" als je denkt dat het verschillende rasters zijn. Voeg niets anders toe.} \\
            }
        }
    }
    & \\ \\

    \theutterance \stepcounter{utterance}  
    & & & \multicolumn{4}{p{0.6\linewidth}}{
        \cellcolor[rgb]{0.9,0.9,0.9}{
            \makecell[{{p{\linewidth}}}]{
                \texttt{\tiny{[GM$\langle$P2]}}
                \texttt{BESLISSING: verschillende rasters} \\
            }
        }
    }
     \\ \\

    \theutterance \stepcounter{utterance}  
    & & & \multicolumn{2}{p{0.3\linewidth}}{
        \cellcolor[rgb]{0.9,0.9,0.9}{
            \makecell[{{p{\linewidth}}}]{
                \texttt{\tiny{[GM$|$GM]}}
                \texttt{continue} \\
            }
        }
    }
    & & \\ \\

    \theutterance \stepcounter{utterance}  
    & & & \multicolumn{2}{p{0.3\linewidth}}{
        \cellcolor[rgb]{0.9,0.9,0.9}{
            \makecell[{{p{\linewidth}}}]{
                \texttt{\tiny{[GM$|$GM]}}
                \texttt{success} \\
            }
        }
    }
    & & \\ \\

    \theutterance \stepcounter{utterance}  
    & & \multicolumn{4}{p{0.6\linewidth}}{
        \cellcolor[rgb]{0.9,0.9,0.9}{
            \makecell[{{p{\linewidth}}}]{
                \texttt{\tiny{[P1$\langle$GM]}}
                \texttt{ANTWOORD: In mijn raster is het eerste vakje van rij 3 (rij 3, kolom 1) leeg (▢). Dat is hetzelfde als bij jou. Maar de rasters zijn nog steeds verschillend, zoals eerder vastgesteld: bij mij is rij 1, kolom 1 gevuld (bij jou leeg) en bij jou is rij 2, kolom 5 leeg (bij mij gevuld).} \\
\texttt{Neem nu een beslissing. Wat denk je: zijn jouw raster en het andere beschreven raster hetzelfde? Schrijf "BESLISSING: hetzelfde raster" als je denkt dat het hetzelfde raster is of "BESLISSING: verschillende rasters" als je denkt dat het verschillende rasters zijn. Voeg niets anders toe.} \\
            }
        }
    }
    & \\ \\

    \theutterance \stepcounter{utterance}  
    & \multicolumn{4}{p{0.6\linewidth}}{
        \cellcolor[rgb]{0.9,0.9,0.9}{
            \makecell[{{p{\linewidth}}}]{
                \texttt{\tiny{[P1$\rangle$GM]}}
                \texttt{BESLISSING: verschillende rasters} \\
            }
        }
    }
    & & \\ \\

    \theutterance \stepcounter{utterance}  
    & & & \multicolumn{2}{p{0.3\linewidth}}{
        \cellcolor[rgb]{0.9,0.9,0.9}{
            \makecell[{{p{\linewidth}}}]{
                \texttt{\tiny{[GM$|$GM]}}
                \texttt{continue} \\
            }
        }
    }
    & & \\ \\

    \theutterance \stepcounter{utterance}  
    & & & \multicolumn{2}{p{0.3\linewidth}}{
        \cellcolor[rgb]{0.9,0.9,0.9}{
            \makecell[{{p{\linewidth}}}]{
                \texttt{\tiny{[GM$|$GM]}}
                \texttt{success} \\
            }
        }
    }
    & & \\ \\

\end{supertabular}
}

\end{document}
